% !TeX root = ../tesis.tex
\chapter{Propiedades ópticas}
\label{chapter:Resultados}


Una vez establecidos los modelos geométricos de los eritrocitos sanos y enfermos, sus funciones dieléctricas y ajustado los parámetros numéricos para la convergencia del método FIT, se calculan sus propiedades ópticas. Este análisis permite evaluar cómo las variaciones simultáneas en la morfología celular y en la concentración de hemoglobina modifican la respuesta óptica de los eritrocitos e identificar regiones espectrales y angulares en las que dichas diferencias se intensifican. En este capítulo se presentan los resultados obtenidos para el eritrocito sano y para dos alteraciones morfológicas representativas: el esferocito y el microcito. Las simulaciones se realizaron en el intervalo espectral de 400 a 600 nm. Para describir las propiedades ópticas de los eritrocitos se emplean las funciones dieléctricas complejas obtenidas en la Sección~\ref{section:optical_properties}, las cuales dependen tanto de la longitud de onda como de la concentración de hemoglobina corpuscular media (CHCM), mientras que el plasma se considera con una permitividad constante en todo el intervalo espectral. En primer lugar, se evalúa la capacidad del modelo numérico para reproducir resultados reportados en la literatura mediante la simulación de un eritrocito descrito con el modelo de Fung \cite{fungBiomechanicsMechanicalProperties1990} a una longitud de onda de $632.8~\text{nm}$. Una vez validado el modelo, se presentan los resultados correspondientes al eritrocito sano modelado mediante un óvalo de Cassini. Posteriormente, se comparan las propiedades ópticas del eritrocito sano, el esferocito y el microcito en el rango de interés por medio sus eficiencias de extinción, así como de los patrones de esparcimiento en tres longitudes de onda representativas de las bandas de absorción de la hemoglobina. Posteriormente, se estudia la distribución angular del esparcimiento mediante el factor de anisotropía, las distribuciones angulares y las contribuciones integradas en las regiones frontal, lateral y de retroesparcimiento. A partir de este análisis se identifican regiones angulares en las que el contraste entre eritrocitos sanos y enfermos aumenta.

\section{Contraste con la literatura}
\label{section:Res_literatura}

Con el fin de evaluar la capacidad del modelo para reproducir resultados reportados en la literatura, se compararon las simulaciones realizadas en este trabajo con los resultados de Tsinopoulos et al. \cite{tsinopoulosLightScatteringAggregated2002}, quienes emplearon el modelo de eritrocito propuesto por Fung \cite{fungBiomechanicsMechanicalProperties1990}, descrito por la ecuación en coordenadas cilíndricas
%
\begin{equation}
	z(\rho)=\left[1-\left(\dfrac{\rho}{a}\right)^2\right]^{1/2}\left[0.72+4.152\left(\dfrac{\rho}{a}\right)^2-3.426\left(\dfrac{\rho}{a}\right)^4\right],
	\label{eq:Tsang_eri}
\end{equation}
%
donde $\rho^2=x^2+y^2$ y $a$ es el radio del eritrocito, para el cual se empleó el valor de $3.825\,\mu$m. Esta geometría corresponde a un eritrocito con un grosor mínimo de $1.44\,\mu$m y un grosor máximo de $2.84\,\mu$m, un volumen de $97.91~\mu\text{m}^3$ y un área de $129.95~\mu\text{m}^2$. Cabe señalar que esta parametrización fue desarrollada para describir eritrocitos sanos, por lo que su aplicación se restringe a este tipo de células. En sus simulaciones, Tsinopoulos et al. emplearon el \textit{Boundary Element Method} (BEM)
y consideraron un eritrocito orientado con su eje de simetría a lo largo de la dirección $z$, iluminado por una onda plana propagándose en la misma dirección y polarizada en $x$, con una longitud de onda de $632.8~\text{nm}$. En esta longitud de onda consideraron un índice de refracción de 1.402 para el eritrocito y de 1.335 para el plasma, lo que corresponde a un índice de refracción relativo de aproximadamente 1.05. Estas mismas condiciones fueron reproducidas en CST empleando los parámetros de convergencia determinados en la Sección \ref{section:convergence}. En la Fig. \ref{fig:Tsinopoulos} se muestran los resultados de la sección transversal diferencial de esparcimiento $\dd{C_{\text{sca}}/\dd{\Omega}}$ [Ec. \eqref{C_sca_general}] en función del ángulo de esparcimiento $\theta$  obtenidos mediante CST con puntos azules y los obtenidos mediante BEM extraídos de la Ref.~\cite{tsinopoulosLightScatteringAggregated2002}. En la Fig. \ref{subfig:Tsi_par} se muestra el corte de $\varphi=0^{\circ}$ (plano paralelo a la dirección de polarización de la onda incidente), mientras que en la Fig. \ref{subfig:Tsi_perp} se muestra el corte de $\varphi=90^{\circ}$ (plano perpendicular a la dirección de polarización de la onda incidente). %
%
\begin{figure}[h!]
	\centering
	\sidesubfloat[First image]{\hspace{-0.1cm}{
			\includegraphics[width=0.45\textwidth]{4-Resultados/comparacionpaper0.pdf}\label{subfig:Tsi_par}}}\hspace{0.1cm}
	\sidesubfloat[Second image]{\hspace{-0.1cm}{	\includegraphics[width=0.45\textwidth]{4-Resultados/comparacionpaper90.pdf}\label{subfig:Tsi_perp}}}
	
	
	\caption{Sección transversal diferencial de esparcimiento $\dd{C_{\text{sca}}/\dd{\Omega}}$ en función del ángulo de esparcimiento $\theta$ para un eritrocito con geometría descrita por Fung \cite{fungBiomechanicsMechanicalProperties1990} y con radio igual a 3.825 $\mu$m. El eritrocito fue iluminado por una onda plana propagándose en la dirección $z$ y polarizada en $x$, con una longitud de onda de $632.8~\text{nm}$. Se consideró un índice de refracción relativo de 1.05. Los puntos azules muestran los resultados obtenidos mediante el método FIT y los naranjas mediante BEM en la Ref.  \cite{tsinopoulosLightScatteringAggregated2002}}
	\label{fig:Tsinopoulos}
\end{figure}
%

Se observa que ambos métodos reproducen la misma estructura angular general del patrón de esparcimiento. En particular, en el plano paralelo  ($\varphi=0^{\circ}$), la diferencia entre en la posición angular de los máximos locales obtenida en la Ref.~\cite{tsinopoulosLightScatteringAggregated2002} y la obtenida en CST es de $0.63^\circ$. De manera similar, en el plano perpendicular ($\varphi=90^{\circ}$), la diferencia promedio en la localización de los máximos locales es de  $0.46^\circ$. Por otro lado, se encuentran diferencias en la amplitud de $\dd{C_{\text{sca}}/\dd{\Omega}}$ con un error relativo promedio de 0.85 y 0.54 para $\varphi=0^{\circ}$ y $\varphi=90^{\circ}$, respectivamente.  Este tipo de diferencias tanto en la amplitud como en la posición de máximos obtenidos mediante diferentes métodos numéricos han sido reportadas previamente en las Refs.~\cite{kupresakBenchmarkingSoftwareTools2017, amirjaniFiniteBoundaryElement}, donde se señala que pueden persistir aún cuando cada método haya sido convergido de manera independiente debido a las diferencias en su formulación numérica y en la discretización espacial. En conjunto, la concordancia observada en la estructura angular y en la localización de los máximos indica que el modelo implementado en CST reproduce el patrón angular de esparcimiento reportado por Tsinopoulos et al., aunque persisten diferencias cuantitativas en su amplitud \cite{kupresakBenchmarkingSoftwareTools2017}. Estos resultados proporcionan una validación adicional de la implementación numérica empleada en este trabajo.
